% Generated draft: requires mathematical and source review.
\documentclass[11pt]{article}
\usepackage[margin=25mm]{geometry}
\usepackage{fontspec}
\setmainfont{Latin Modern Roman}
\usepackage{amsmath,amssymb,mathtools,unicode-math}
\setmathfont{Latin Modern Math}
\usepackage{enumitem,needspace,xurl,hyperref}
\setlength{\parindent}{0pt}
\setlength{\parskip}{5pt}
\setlength{\emergencystretch}{3em}
\newcommand{\N}{\mathbb{N}}
\newcommand{\Z}{\mathbb{Z}}
\newcommand{\Q}{\mathbb{Q}}
\newcommand{\R}{\mathbb{R}}
\newcommand{\C}{\mathbb{C}}
\newcommand{\F}{\mathbb{F}}
\newcommand{\A}{\mathbb{A}}
\newcommand{\PP}{\mathbb{P}}
\newcommand{\E}{\mathbb{E}}
\newcommand{\eps}{\varepsilon}
\newcommand{\dd}{\,\mathrm{d}}
\begin{document}
\section*{UnsolvedMath ID: 800076}
\textbf{Source problem code:} catalogue entry
\subsection{Definitions and mathematical statement}
Let
\[
D^2=\{x\in\R^2:\|x\|\leq 1\},\qquad *=0,
\]
and, following the source's explicit description, let
\[
G=\{f:D^2\longrightarrow D^2:\ f\text{ is a smooth orientation-preserving diffeomorphism and }f(0)=0\},
\]
with group operation composition. No requirement that the boundary be fixed pointwise is stated. The source denotes this group by \(\mathrm{Diff}_0(D^2,*)\); it does not separately explain the subscript \(0\).

A real-valued quasimorphism on \(G\) is a function \(q:G\to\R\) with finite defect
\[
D(q)=\sup_{g,h\in G}|q(gh)-q(g)-q(h)|<\infty.
\]
It is homogeneous if
\[
q(g^n)=nq(g)\qquad(g\in G,\ n\in\Z).
\]
Write \(\mathcal Q_{\mathrm{hom}}(G)\) for the real vector space of all such homogeneous quasimorphisms, with pointwise operations. This definition imposes no continuity or measurability condition and takes no quotient by homomorphisms or bounded functions; the source does not explicitly discuss these conventions.

The source asserts that there are two bounded Euler classes associated with this group, one induced by restriction to \(\partial D^2\), and the other induced by the derivative at \(0\). It further asserts that their difference gives a nontrivial quasimorphism on \(G\). These are supplied background assertions; the record does not specify the class normalizations or the construction of that quasimorphism.

\textbf{Question 17 (Nariman, Yazdi, 2025).} Is
\[
\dim_{\R}\mathcal Q_{\mathrm{hom}}(G)=1?
\]
This is the question under the explicit group description and quasimorphism conventions above. Any intended component restriction, regularity requirement, or quotient convention requires clarification from the source.
\subsection{Short English statement}
For the group of smooth orientation-preserving diffeomorphisms of the closed disk fixing its marked origin, is the real vector space of homogeneous quasimorphisms one-dimensional? The source supplies a nontrivial quasimorphism arising from two bounded Euler classes. The precise conventions for the group notation and the space of quasimorphisms require review.
\subsection{Sources}
\begin{itemize}
\item Local UnsolvedMath registry, permanent ID \(800076\), exact source problem code \texttt{catalogue entry}, supplied complete record titled \emph{Foliations and Diffeomorphism Groups --- Question 17}. Upstream revision: \texttt{d371ee917ac9647fbb06cd6fc529deec6f3dda82}.
\item Nariman, Yazdi, \emph{Problem List on Foliations and Diffeomorphism Groups} (2025), Question 17. Supplied source URL: \url{https://arxiv.org/abs/2503.18274}. The record attributes the question to Nariman and Yazdi, with proposed year 2025.
\item Supplied provenance records access on \texttt{2026-10-16} and extraction by \texttt{arxiv\_latex}. This access information is supplied provenance, not an independently verified source visit. The registry labels the problem \emph{open}; that label is supplied curation, not an independently verified literature assessment.
\end{itemize}
\end{document}
